\section{Conclusion}
\label{sec:conclusion}

This paper has developed a stylized two-region framework for evaluating public co-creation hubs in thin regional ecosystems. The model separates two policy margins that are often conflated in practice: \emph{local coordination} inside the target region and \emph{cross-regional renewal} through outside linkages. Its principal distortion is a dynamic coordination externality. Decentralized actors can value current coordination and directly contracted outside links, but they do not fully internalize how current local interaction changes the future novelty available to the regional system.

Three results follow. First, decentralized support and planner evaluation need not coincide. The same hub can be under-supported when renewal gains are neglected or over-supported when visible local activity masks future redundancy losses. Second, a hub is socially efficient as a stand-alone intervention only under limited conditions: its coordination gain must be large enough to offset fixed cost and the dynamic erosion of novelty generated by repeated local interaction. Third, pipeline policy has positive stand-alone value, so the relevant policy comparison is often not hub versus no hub, but pipeline alone versus a bundle that adds a hub to a renewal strategy. The bundle is preferred only under specific parameter conditions.

These conclusions imply a more demanding regional policy evaluation. A hub should not be judged only by attendance, occupancy, meetings, or other short-run utilization measures. A theory-consistent evaluation should also ask whether new ties continue to form, whether repeated collaborator ratios rise over time, whether outside-region collaboration expands under pipeline policy, and whether a bundled strategy outperforms a pipeline-only baseline. In regional science terms, the key issue is not merely whether a local venue is active, but whether the regional system combines internal coordination with enough external renewal to remain dynamically productive in the model's minimal two-period sense.

For that reason, the most informative empirical setting is not a simple treated-versus-untreated comparison, but a three-way comparison among no organized intervention, pipeline-oriented renewal without a hub, and the bundled strategy. Only that design can show whether a hub adds value beyond an already active renewal policy.

The framework is intentionally parsimonious, and that parsimony defines its limits. Matching rates, outside-link productivity, and the outside pool are all reduced-form objects. Network formation is not endogenous, and neither migration nor location choice is modeled. The analysis is therefore best interpreted as a conditional policy theory rather than as a full structural account of regional evolution. Its purpose is to isolate one principal welfare wedge and show how that wedge changes the evaluation of hub policy in thin regional ecosystems. The partial-internalization extension also makes clear that the benchmark is not all-or-nothing: as decentralized actors internalize more of the continuation value, the support-threshold gap shrinks and vanishes in the full-internalization limit.

Within those limits, the paper offers a conditional implication for the design of place-based policy. In thin regional ecosystems, public hubs are most plausibly justified not as stand-alone venues that generate visible activity, but as elements of a broader regional strategy that combines local coordination with cross-regional renewal. Whether such a bundle is desirable remains conditional, but the relevant benchmark for judging hub policy changes once renewal is treated as a policy margin in its own right.
